\chapter{Gaussian and Eisenstein Integers}

\begin{center}\itshape
In $\Z[i]$, the prime $5$ splits as $(2+i)(2-i)$, the prime $3$ stays prime, and $2=-i(1+i)^2$.\\
This one sentence contains the two-square theorem, and the same game with $\omega=e^{2i\pi/3}$ proves Fermat for cubes.
\end{center}

\section{Three rings and their norms}

Let $\omega=e^{2i\pi/3}=\frac{-1+i\sqrt3}2$, so $\omega^2+\omega+1=0$ and $\bar\omega=\omega^2$. We study the rings
\[
\Z[i]=\{a+bi\},\qquad\Z[\omega]=\{a+b\omega\},\qquad\Z[\sqrt{-2}]=\{a+b\sqrt{-2}\}\qquad(a,b\in\Z),
\]
subrings of $\C$ (closed under products because $i^2=-1$, $\omega^2=-1-\omega$, $\sqrt{-2}^2=-2$) and stable under complex conjugation. Their \emph{norm} is $N(z)=z\bar z=|z|^2$, a non-negative integer:
\[
N(a+bi)=a^2+b^2,\qquad N(a+b\omega)=a^2-ab+b^2,\qquad N(a+b\sqrt{-2})=a^2+2b^2,
\]
the middle one because $|a+b\omega|^2=\left(a-\frac b2\right)^2+\frac34b^2$. The norm is multiplicative and vanishes only at $0$.

In such a ring $R$, a \emph{unit} is an invertible element; $a,b$ are \emph{associates} if $a=ub$ with $u$ a unit; and a non-zero non-unit $\pi$ is \emph{irreducible} if $\pi=ab$ forces $a$ or $b$ to be a unit.

\begin{proposition}\label{prop:quadunits}
In each of the three rings, $z$ is a unit if and only if $N(z)=1$. The units are $\pm1,\pm i$ in $\Z[i]$, the six numbers $\pm1,\pm\omega,\pm\omega^2$ in $\Z[\omega]$, and $\pm1$ in $\Z[\sqrt{-2}]$. An element whose norm is a prime number is irreducible.
\end{proposition}
\begin{proof}
If $zw=1$ then $N(z)N(w)=1$ in $\N$, so $N(z)=1$. Conversely if $z\bar z=1$ then $\bar z\in R$ is the inverse. Solving $N=1$: $a^2+b^2=1$ gives the four Gaussian units; $a^2+2b^2=1$ gives $\pm1$; and $4(a^2-ab+b^2)=(2a-b)^2+3b^2=4$ gives $b=0$, $a=\pm1$, or $b=\pm1$, $2a-b=\pm1$, that is $1,-1,\omega,-\omega,-1-\omega=\omega^2,1+\omega=-\omega^2$. Finally if $N(\pi)$ is prime and $\pi=ab$, then $N(a)N(b)$ is prime, so $N(a)=1$ or $N(b)=1$.
\end{proof}

\begin{think}
To divide $11+3i$ by $1+8i$ in $\Z[i]$, compute the quotient $\frac{11+3i}{1+8i}$ in $\C$ and round it to a nearby Gaussian integer. How far can a complex number be from the nearest point of the square lattice $\Z[i]$? Of the triangular lattice $\Z[\omega]$? Of the lattice $\Z[\sqrt{-5}]$, whose cells are $1\times\sqrt5$ rectangles? Which of these distances are less than $1$, and why does that matter?
\end{think}

\begin{proposition}[Euclidean division]\label{prop:eucl}
In each of the three rings, for $a,b$ with $b\ne0$, there are $q,r$ with $a=bq+r$ and $N(r)<N(b)$.
\end{proposition}
\begin{proof}
Let $z=\frac ab\in\C$. In each case $z=s+t\theta$ with $s,t\in\Q$ and $\theta\in\{i,\omega,\sqrt{-2}\}$: multiply numerator and denominator by $\bar b$, and note $N(b)\in\Z$, $a\bar b\in R$. Choose integers $m,n$ with $|s-m|\le\frac12$, $|t-n|\le\frac12$, and let $q=m+n\theta$, $\sigma=s-m$, $\tau=t-n$. Then $\left|\frac ab-q\right|^2=|\sigma+\tau\theta|^2$ is $\sigma^2+\tau^2\le\frac12$, or $\sigma^2-\sigma\tau+\tau^2\le\frac34$, or $\sigma^2+2\tau^2\le\frac34$. In all cases it is $<1$, and $r=a-bq\in R$ has $N(r)=N(b)\left|\frac ab-q\right|^2<N(b)$.
\end{proof}

\section{Unique factorisation}

\begin{theorem}\label{thm:ufd}
Let $R$ be one of the three rings (or any integral domain with a multiplicative norm $N:R\to\N$ satisfying Propositions~\ref{prop:quadunits} and~\ref{prop:eucl}).
\begin{enumerate}[label=(\alph*)]
\item (B\'ezout) For $a,b\in R$ not both zero, there is $d\in R$ with $aR+bR=dR$; $d$ divides $a$ and $b$, and every common divisor divides $d$.
\item (Gauss, Euclid) An irreducible $\pi$ is \emph{prime}: if $\pi\mid ab$ then $\pi\mid a$ or $\pi\mid b$.
\item Every non-zero non-unit is a product of irreducibles, uniquely up to order and to associates.
\end{enumerate}
\end{theorem}
\begin{proof}
(a) $I=aR+bR$ is stable under sums and under multiplication by elements of $R$. Let $d\in I$ be non-zero of least norm. For $x\in I$, divide: $x=dq+r$ with $r=x-dq\in I$ and $N(r)<N(d)$, so $r=0$. Hence $I=dR$. Then $a,b\in dR$, and $d=au+bv$ is divisible by every common divisor.

(b) Suppose $\pi\mid ab$ and $\pi\nmid a$. Let $\pi R+aR=dR$. Then $d\mid\pi$, so $d$ is a unit or an associate of $\pi$; in the second case $\pi\mid d\mid a$, excluded. So $d$ is a unit and $1=\pi x+ay$ for some $x,y$. Then $b=\pi xb+aby$ is divisible by $\pi$.

(c) Existence, by strong induction on $N(z)\ge2$: if $z$ is not irreducible, $z=ab$ with $a,b$ non-units, so $N(a),N(b)\ge2$ and $N(a),N(b)<N(z)$. Uniqueness, by induction on the number of factors: if $\pi_1\cdots\pi_r=u\,\rho_1\cdots\rho_s$, then $\pi_1$ divides some $\rho_j$ by (b), so $\rho_j=v\pi_1$ with $v$ a unit (as $\rho_j$ is irreducible); cancel $\pi_1$ and continue.
\end{proof}

As in $\Z$, we may then speak of the \emph{valuation} $v_\pi(z)$ of $z$ at an irreducible $\pi$, and two consequences used repeatedly follow exactly as in $\Z$: if $\gcd(a,b)=1$ (no common irreducible factor) and $ab=uc^k$ with $u$ a unit, then $a$ and $b$ are each a unit times a $k$-th power; and integers coprime in $\Z$ are coprime in $R$ (a B\'ezout relation in $\Z$ is one in $R$).

\begin{remark}
Unique factorisation fails in general: in $\Z[\sqrt{-5}]$, $6=2\cdot3=(1+\sqrt{-5})(1-\sqrt{-5})$ (Exercise~\ref{ex:zsqrt5}). Kummer's ideal numbers, the subject of Chapter~11, were invented to repair this.
\end{remark}

\section{Gaussian primes and sums of two squares}

\begin{think}
Which of the primes $2,3,5,7,11,13$ factor in $\Z[i]$? Try to write each one as $(a+bi)(a-bi)$. Which ones refuse, and what do they have in common? Compare with Chapter~7.
\end{think}

\begin{theorem}\label{thm:gaussprimes}
Up to associates, the irreducibles of $\Z[i]$ are:
\begin{enumerate}[label=(\roman*)]
\item $1+i$, with $2=-i(1+i)^2$;
\item the primes $q\equiv3\pmod4$ of $\Z$;
\item for each prime $p\equiv1\pmod4$, two non-associate irreducibles $\pi=a+bi$ and $\bar\pi=a-bi$ with $p=\pi\bar\pi=a^2+b^2$.
\end{enumerate}
\end{theorem}
\begin{proof}
$N(1+i)=2$ and $N(a+bi)=p$ for $p\equiv1\pmod4$ (Theorem~\ref{thm:twosq}), so these are irreducible by Proposition~\ref{prop:quadunits}. For $\pi=a+bi$ with $a^2+b^2=p$ odd, $\bar\pi$ is not an associate of $\pi$: otherwise $\pi$ would divide $\pi+\bar\pi=2a$ and $\pi-\bar\pi=2bi$, so $p=N(\pi)$ would divide $4a^2$ and $4b^2$, hence $p\mid a$, $p\mid b$, and $p^2\mid a^2+b^2=p$. A prime $q\equiv3\pmod4$ is irreducible: if $q=\alpha\beta$ with non-units, then $N(\alpha)N(\beta)=q^2$ forces $N(\alpha)=q$, but $q$ is not a sum of two squares.
Conversely let $\pi$ be irreducible. It divides $N(\pi)=\pi\bar\pi$, a product of prime numbers, hence (Theorem~\ref{thm:ufd}(b)) divides one prime number $p$. Then $N(\pi)\mid N(p)=p^2$. Factoring $p$ by (i)--(iii) and using uniqueness, $\pi$ is an associate of one of the listed factors of $p$.
\end{proof}

Figure~\ref{fig:gausseis} draws the irreducibles of both rings. The pattern has the symmetry of the units: a rotation by $90^\circ$ in $\Z[i]$, and by $60^\circ$ in $\Z[\omega]$ (Theorem~\ref{thm:eisprimes} below).

\begin{figure}[ht]
\centering
\fig{ch10-primes}
\caption{The irreducibles of norm at most $225$ in $\Z[i]$ (left) and $\Z[\omega]$ (right). Blue: the factors $\pi$ of a prime $p=\pi\bar\pi$ with $p\equiv1\pmod4$ (left) or $p\equiv1\pmod3$ (right). Red: the primes $q\equiv3\pmod4$ (left) or $q\equiv2\pmod3$ (right), which stay irreducible, with their associates. Green: the associates of $1+i$ and of $1-\omega$, whose squares are $2$ and $3$ up to units.}
\label{fig:gausseis}
\end{figure}

We can now prove the formula for $r_2(n)=\#\{(x,y)\in\Z^2:x^2+y^2=n\}$ announced in Chapter~7.

\begin{theorem}[Jacobi]\label{thm:r2}
For $n\ge1$, $r_2(n)=4\displaystyle\sum_{d\mid n}\chi_4(d)$, where $\chi_4(d)$ is $0$ for even $d$, $1$ for $d\equiv1$ and $-1$ for $d\equiv3\pmod4$.
\end{theorem}
\begin{proof}
$r_2(n)$ is the number of $z\in\Z[i]$ with $N(z)=n$. Write $n=2^g\prod_jp_j^{e_j}\prod_kq_k^{f_k}$ with $p_j\equiv1$, $q_k\equiv3\pmod4$, and fix $\pi_j$ with $p_j=\pi_j\bar\pi_j$. By Theorems~\ref{thm:ufd} and~\ref{thm:gaussprimes}, every non-zero $z$ is uniquely
\[
z=u\,(1+i)^{g'}\prod_j\pi_j^{a_j}\bar\pi_j^{b_j}\prod_kq_k^{c_k}\cdots,\qquad N(z)=2^{g'}\prod_jp_j^{a_j+b_j}\prod_kq_k^{2c_k}\cdots
\]
with $u$ one of $4$ units. So $N(z)=n$ requires $g'=g$, $a_j+b_j=e_j$ ($e_j+1$ choices), $2c_k=f_k$, and no other primes. Hence $r_2(n)=4\prod_j(e_j+1)$ if all $f_k$ are even, and $0$ otherwise.
On the other side, $\chi_4$ is completely multiplicative, so $\delta(n)=\sum_{d\mid n}\chi_4(d)$ is multiplicative (Proposition~\ref{prop:conv}), with $\delta(2^g)=1$, $\delta(p^e)=e+1$, and $\delta(q^f)=1-1+\dots\pm1$, which is $1$ if $f$ is even and $0$ if $f$ is odd. The two expressions agree.
\end{proof}

\section{Eisenstein integers and Fermat for cubes}

Let $\lambda=1-\omega$. Then $N(\lambda)=1+1+1=3$, so $\lambda$ is irreducible, and $\lambda^2=1-2\omega+\omega^2=-3\omega$, so $3=-\omega^2\lambda^2$.

\begin{theorem}\label{thm:eisprimes}
Up to associates, the irreducibles of $\Z[\omega]$ are $\lambda=1-\omega$; the primes $q\equiv2\pmod3$; and, for each prime $p\equiv1\pmod3$, two non-associate irreducibles $\pi,\bar\pi$ with $p=\pi\bar\pi=N(\pi)$. Moreover every $z\in\Z[\omega]$ is congruent to $0$, $1$ or $-1$ modulo $\lambda$.
\end{theorem}
\begin{proof}
Let $p\equiv1\pmod3$. By Exercise~\ref{ex:minus3}, $a^2\equiv-3\pmod p$ for some $a$, so $p\mid a^2+3=(a+\sqrt{-3})(a-\sqrt{-3})$, where $\sqrt{-3}=1+2\omega\in\Z[\omega]$. But $p$ divides neither factor, since $\frac{a\pm(1+2\omega)}p$ has coefficient $\pm\frac2p\notin\Z$ on $\omega$. So $p$ is not prime, hence not irreducible (Theorem~\ref{thm:ufd}(b)): $p=\pi\alpha$ with non-units, and $N(\pi)N(\alpha)=p^2$ gives $N(\pi)=p=\pi\bar\pi$. If $\bar\pi$ were an associate of $\pi=a+b\omega$, then $\pi$ would divide $\pi+\bar\pi=2a-b$ and $\pi-\bar\pi=b(\omega-\omega^2)$, of norm $3b^2$; so $p=N(\pi)$ would divide $(2a-b)^2$ and $3b^2$, hence $p\mid b$ (as $p\ne3$), then $p\mid2a$ and $p\mid a$, and $p^2\mid N(\pi)=p$, absurd. For $q\equiv2\pmod3$: a factor of norm $q$ would give $4q=(2a-b)^2+3b^2\equiv(2a-b)^2\in\{0,1\}\pmod3$, while $4q\equiv2$; so $q$ is irreducible. As for Gaussian integers, every irreducible divides a prime number and is therefore in the list (with $3=-\omega^2\lambda^2$).
Finally $\omega\equiv1\pmod\lambda$, so $a+b\omega\equiv a+b\pmod\lambda$; and $\lambda\mid3$, so an integer is congruent modulo $\lambda$ to its residue in $\{-1,0,1\}$ modulo $3$.
\end{proof}

\begin{lemma}\label{lem:cube}
If $\lambda\nmid\theta$, then $\theta^3\equiv\pm1\pmod{\lambda^4}$.
\end{lemma}
\begin{proof}
By Theorem~\ref{thm:eisprimes}, $\pm\theta=1+\beta\lambda$ for a sign and some $\beta$. Since $1-\omega^2=(1-\omega)(1+\omega)=-\omega^2\lambda$,
\[
(\pm\theta)^3-1=(\pm\theta-1)(\pm\theta-\omega)(\pm\theta-\omega^2)=\beta\lambda\,(\beta\lambda+\lambda)\,(\beta\lambda-\omega^2\lambda)=\lambda^3\beta(\beta+1)(\beta-\omega^2).
\]
Modulo $\lambda$, $\omega^2\equiv1$, so $\beta,\beta+1,\beta-\omega^2$ are congruent to $\beta,\beta+1,\beta-1$, and one of these is $\equiv0$ since $\beta\equiv0,\pm1$. So $\lambda^4\mid(\pm\theta)^3-1$.
\end{proof}

\begin{think}
Compute the cubes modulo $9$. Which values occur for cubes of integers prime to $3$? Deduce that if $x^3+y^3=z^3$ with integers, then $3$ divides one of $x,y,z$. The proof below makes the same idea work with $\lambda=1-\omega$ instead of $3$.
\end{think}

\begin{theorem}[Euler; Gauss]\label{thm:flt3}
The equation $x^3+y^3=z^3$ has no solution in non-zero integers.
\end{theorem}
\begin{proof}
Suppose $x^3+y^3=z^3$ with $xyz\ne0$; dividing by the gcd we may assume $x,y,z$ pairwise coprime in $\Z$ (a prime dividing two of them divides the third). If $\lambda$ divides an integer $m$, then $3=N(\lambda)\mid N(m)=m^2$, so $3\mid m$; hence $\lambda$ divides at most one of $x,y,z$.

\emph{Case 1: $\lambda\nmid xyz$.} By Lemma~\ref{lem:cube}, $0=x^3+y^3-z^3\equiv\pm1\pm1\pm1\pmod{\lambda^4}$. The possible values are $\pm1$ and $\pm3$; neither is divisible by $\lambda^4$, since $\lambda\nmid1$ and $v_\lambda(3)=2$. Contradiction.

\emph{Case 2: $\lambda$ divides one of them.} Renaming (and replacing $z$ by $-z$, as $(-z)^3=-z^3$), we get $\alpha^3+\beta^3=\varepsilon\lambda^{3n}\gamma^3$ with $\varepsilon$ a unit, $n\ge1$, and $\alpha,\beta,\gamma$ non-zero, prime to $\lambda$, with $\gcd(\alpha,\beta)=1$. We show that no such equation holds in $\Z[\omega]$, by descent on $n$.

(i) \emph{$n\ge2$.} By Lemma~\ref{lem:cube}, $\alpha^3+\beta^3\equiv\pm1\pm1\in\{0,\pm2\}\pmod{\lambda^4}$, and it is divisible by $\lambda^3$, while $\lambda\nmid2$ (as $3\nmid N(2)=4$). So $\lambda^4\mid\alpha^3+\beta^3$, i.e.\ $3n\ge4$.

(ii) \emph{Factor.} $(\alpha+\beta)(\alpha+\omega\beta)(\alpha+\omega^2\beta)=\varepsilon\lambda^{3n}\gamma^3$. The three factors differ pairwise by $\lambda\beta$ times a unit (e.g.\ $(\alpha+\beta)-(\alpha+\omega\beta)=\lambda\beta$, $(\alpha+\omega\beta)-(\alpha+\omega^2\beta)=\omega\lambda\beta$), so they are all congruent modulo $\lambda$; as their product is divisible by $\lambda$, all three are. And no two are divisible by $\lambda^2$, since their difference $\lambda\beta\cdot$unit is not. So two of them have $v_\lambda=1$ and the third has $v_\lambda=3n-2\ge4$. Replacing $\beta$ by $\omega\beta$ or $\omega^2\beta$ (which does not change $\beta^3$) permutes the factors, so we may assume the third is $\alpha+\beta$. A common divisor of two factors divides $\lambda\beta$ and also $\lambda\alpha$ (e.g.\ $\omega(\alpha+\beta)-(\alpha+\omega\beta)=-\lambda\alpha$), hence divides $\lambda$. So
\[
\alpha+\beta=\lambda^{3n-2}\kappa_1,\qquad\alpha+\omega\beta=\lambda\kappa_2,\qquad\alpha+\omega^2\beta=\lambda\kappa_3,
\]
with $\kappa_1,\kappa_2,\kappa_3$ pairwise coprime, prime to $\lambda$, and $\kappa_1\kappa_2\kappa_3=\varepsilon\gamma^3$. By unique factorisation, $\kappa_j=\varepsilon_j\theta_j^3$ with units $\varepsilon_j$ and $\theta_j$ prime to $\lambda$.

(iii) \emph{Descend.} Since $1+\omega+\omega^2=0$,
$(\alpha+\beta)+\omega(\alpha+\omega\beta)+\omega^2(\alpha+\omega^2\beta)=(1+\omega+\omega^2)\alpha+(1+\omega^2+\omega^4)\beta=0$. Substituting and dividing by $\omega\varepsilon_2\lambda$:
\[
\theta_2^3+\eta\,\theta_3^3=\eta'\lambda^{3(n-1)}\theta_1^3
\]
for some units $\eta,\eta'$. As $n\ge2$, the right side is divisible by $\lambda^3$, and by Lemma~\ref{lem:cube} the left side is $\equiv\pm1\pm\eta\pmod{\lambda^4}$, so $\lambda^2\mid\pm1\pm\eta$. Among the units, $\pm1\pm\eta\equiv0\pmod{\lambda^2}$ only for $\eta=\pm1$: indeed $1-\omega=\lambda$ and $1-\omega^2=-\omega^2\lambda$ have $v_\lambda=1$, and $1+\omega=-\omega^2$, $1+\omega^2=-\omega$ are units. So $\eta=\pm1$, and $\theta_2^3+(\pm\theta_3)^3=\eta'\lambda^{3(n-1)}\theta_1^3$ is an equation of the same type with $n-1$ in place of $n$ ($\theta_2,\pm\theta_3$ are coprime, as $\kappa_2,\kappa_3$ are). So, if we choose at the start an equation of this type with $n$ minimal, then $n\ge2$ by (i), and (iii) produces one with exponent $n-1\ge1$: a contradiction. So Case~2 is impossible too.
\end{proof}

\section{Exercises}

\begin{exercise}[$\star$]\label{ex:gaussfactor}
Factor $65$ and $2026$ into Gaussian primes, and compute $\gcd(11+3i,1+8i)$ with Euclid's algorithm in $\Z[i]$.
\end{exercise}

\begin{exercise}[$\star$]\label{ex:zsqrt5}
In $\Z[\sqrt{-5}]$, with $N(a+b\sqrt{-5})=a^2+5b^2$, show that $2$, $3$, $1\pm\sqrt{-5}$ are irreducible and pairwise non-associate. Deduce that $\Z[\sqrt{-5}]$ does not have unique factorisation, and that $2$ is irreducible but not prime there.
\end{exercise}

\begin{exercise}[$\star\star$]\label{ex:gaussgcd}
Let $p\equiv1\pmod4$ be prime and $a^2\equiv-1\pmod p$. Show that $\gcd(p,a+i)$ in $\Z[i]$ is $x+iy$ with $p=x^2+y^2$. Compute it for $p=1013$, $a=45$. (This is the Gaussian form of Cornacchia's algorithm, Chapter~7's notebook.)
\end{exercise}

\begin{exercise}[$\star\star$]\label{ex:euler1000009}
\begin{enumerate}[label=(\alph*)]
\item Show that a prime $p\equiv1\pmod4$ is a sum of two squares in essentially one way: $p=a^2+b^2=c^2+d^2$ implies $\{|a|,|b|\}=\{|c|,|d|\}$.
\item Euler noticed $1000009=1000^2+3^2=972^2+235^2$. Deduce that $1000009$ is composite, and find a factor from $\gcd(1000+3i,972+235i)$.
\end{enumerate}
\end{exercise}

\begin{exercise}[$\star\star$]\label{ex:x3y21}
Show that the only integer solution of $x^3=y^2+1$ is $(1,0)$.
\end{exercise}

\begin{exercise}[$\star\star$]\label{ex:xxyy}
Show that a prime $p$ is of the form $x^2+xy+y^2$ with $x,y\in\Z$ if and only if $p=3$ or $p\equiv1\pmod3$.
\end{exercise}

\begin{exercise}[$\star\star\star$]\label{ex:y2x3m2}
\emph{Fermat's claim.} Show that the only integer solutions of $y^2=x^3-2$ are $(3,\pm5)$. \emph{Hint: factor $y^2+2$ in $\Z[\sqrt{-2}]$.}
\end{exercise}

\begin{beauty}
Theorem~\ref{thm:gaussprimes} says that the way a prime $p$ factors in $\Z[i]$ depends only on $p\bmod4$, and Theorem~\ref{thm:eisprimes} that its factorisation in $\Z[\omega]$ depends only on $p\bmod3$. This is the simplest case of a deep pattern: in a ring generated by roots of unity, the factorisation of $p$ depends only on $p$ modulo $n$, and the general version, class field theory, is one of the summits of twentieth-century mathematics. The same rings carry cubic and biquadratic reciprocity laws, discovered by Gauss and proved by Eisenstein and Jacobi. And Theorem~\ref{thm:flt3} is the first case of Fermat's Last Theorem, proved here by the method Kummer later pushed to all regular primes.
\end{beauty}

\section{Hints}

\begin{hint}{ex:gaussfactor}
$65=5\cdot13$ and $2026=2\cdot1013$, with $1013=22^2+23^2$. For the gcd, divide $11+3i$ by $1+8i$, round the quotient to the nearest Gaussian integer, and continue as in $\Z$.
\end{hint}
\begin{hint}{ex:zsqrt5}
The norm $a^2+5b^2$ never equals $2$ or $3$. So any factor of $2$, $3$ or $1\pm\sqrt{-5}$ (norms $4$, $9$, $6$) must be a unit. Then look at $2\cdot3=(1+\sqrt{-5})(1-\sqrt{-5})$.
\end{hint}
\begin{hint}{ex:gaussgcd}
$p\mid(a+i)(a-i)$ but $p$ divides neither factor. So $p=\pi\bar\pi$, and $\pi$ divides one of the two factors. What is the norm of $\gcd(p,a+i)$?
\end{hint}
\begin{hint}{ex:euler1000009}
(a) $p=(a+bi)(a-bi)=(c+di)(c-di)$, and unique factorisation leaves only units and conjugation as freedom. (b) If $1000009$ were prime, (a) would forbid two essentially different representations.
\end{hint}
\begin{hint}{ex:x3y21}
Factor $y^2+1=(y+i)(y-i)$. Show the factors are coprime (look at $x$ modulo $2$ first), so $y+i$ is a unit times a cube; every unit of $\Z[i]$ is itself a cube.
\end{hint}
\begin{hint}{ex:xxyy}
Replacing $y$ by $-y$, $x^2+xy+y^2$ becomes the norm $x^2-xy+y^2$ of $x+y\omega$. So $p$ has this form if and only if it is a norm, that is, not irreducible in $\Z[\omega]$.
\end{hint}
\begin{hint}{ex:y2x3m2}
$x^3=(y+\sqrt{-2})(y-\sqrt{-2})$ in $\Z[\sqrt{-2}]$, a ring with unique factorisation and units $\pm1$. Show the factors are coprime, write $y+\sqrt{-2}=(a+b\sqrt{-2})^3$, and compare the coefficients of $\sqrt{-2}$.
\end{hint}

\section{Solutions}

\begin{solution}{ex:gaussfactor}
\idea Factor in $\Z$ first, then split each prime with Theorem~\ref{thm:gaussprimes}: $p\equiv1\pmod4$ splits as $(a+bi)(a-bi)$ where $p=a^2+b^2$, $2=-i(1+i)^2$, and $q\equiv3\pmod4$ stays. For the gcd, divide and round to the nearest Gaussian integer, as in the proof of Proposition~\ref{prop:eucl}.
\writeup $65=5\cdot13=(2+i)(2-i)(3+2i)(3-2i)$. $2026=2\cdot1013=-i(1+i)^2(22+23i)(22-23i)$, using $1013=22^2+23^2$ (Exercise~\ref{ex:2026}). For the gcd: $\frac{11+3i}{1+8i}=\frac{(11+3i)(1-8i)}{65}=\frac{35-85i}{65}$, nearest Gaussian integer $1-i$; remainder $11+3i-(1+8i)(1-i)=11+3i-(9+7i)=2-4i$. Then $\frac{1+8i}{2-4i}=\frac{(1+8i)(2+4i)}{20}=\frac{-30+20i}{20}$, round to $-1+i$ (or $-2+i$); remainder $1+8i-(2-4i)(-1+i)=1+8i-(2+6i)=-1+2i$. Then $\frac{2-4i}{-1+2i}=-2$ exactly. So the gcd is $-1+2i$, an associate of $2+i$, of norm $5$: indeed $N(11+3i)=130$ and $N(1+8i)=65$ are both divisible by $5$.
\end{solution}

\begin{solution}{ex:zsqrt5}
\idea To show an element is irreducible, use norms: a non-trivial factor would have a norm strictly between $1$ and the total, and the form $a^2+5b^2$ never takes the values $2$ or $3$. Units have norm $1$, so they are $\pm1$, and associates are just $\pm$.
\writeup As in Proposition~\ref{prop:quadunits}, units have norm $1$, so they are $\pm1$. The norms are $N(2)=4$, $N(3)=9$, $N(1\pm\sqrt{-5})=6$. A non-trivial factorisation of one of them would involve an element of norm $2$ or $3$; but $a^2+5b^2\in\{2,3\}$ is impossible ($b=0$ gives a square, $b\ne0$ gives at least $5$). So all four are irreducible; they are pairwise non-associate because associates are $\pm$ each other. Then $6=2\cdot3=(1+\sqrt{-5})(1-\sqrt{-5})$ are two essentially different factorisations. Finally, $2\mid(1+\sqrt{-5})(1-\sqrt{-5})$ but $2$ divides neither factor ($\frac{1\pm\sqrt{-5}}2\notin\Z[\sqrt{-5}]$), so $2$ is not prime. This shows that the ring has no Euclidean division for any norm-like function: Theorem~\ref{thm:ufd}(b) would fail.
\end{solution}

\begin{solution}{ex:gaussgcd}
\idea $p\mid a^2+1=(a+i)(a-i)$ but $p$ divides neither factor, so $p$ is not a Gaussian prime: $p=\pi\bar\pi$. The irreducible $\pi$ divides $a+i$ or $a-i$; the gcd of $p$ and $a+i$ is then a non-trivial divisor of $p$, which must have norm exactly $p$.
\writeup $p\mid(a+i)(a-i)$ and $p\nmid a\pm i$ (as $\frac{a\pm i}p\notin\Z[i]$). Let $g=\gcd(p,a+i)$. If $g$ were a unit, then by Theorem~\ref{thm:ufd}(a) $1=px+(a+i)y$, and multiplying by $a-i$: $p\mid a-i$, impossible. If $g$ were an associate of $p$, then $p\mid a+i$, impossible. So $g$ is a proper divisor of $p$, and $N(g)$ is a proper divisor of $N(p)=p^2$, so $N(g)=p$: $g=x+iy$ with $x^2+y^2=p$. For $p=1013$, $a=45$ ($45^2+1=2026=2\cdot1013$): $\frac{1013}{45+i}=\frac{1013(45-i)}{2026}=\frac{45-i}2$, round to $22$ (or $23$, or with $\pm\frac i2$ rounded to $0$ or $-i$); remainder $1013-22(45+i)=23-22i$. Then $\frac{45+i}{23-22i}=\frac{(45+i)(23+22i)}{1013}=\frac{1013+1013i}{1013}=1+i$ exactly. So $g=23-22i$ and $1013=23^2+22^2$.
\end{solution}

\begin{solution}{ex:euler1000009}
\idea (a) A representation $p=a^2+b^2$ is a factorisation $p=(a+bi)(a-bi)$ into irreducibles, and unique factorisation says there is essentially only one; the remaining freedom (units, conjugation) only permutes and changes signs. (b) Two essentially different representations of $n$ mean that $n$ factors in two different ways in $\Z[i]$, which cannot happen for a prime; and the gcd of the two Gaussian integers isolates a common factor, whose norm divides $n$.
\writeup (a) $a+bi$ has norm $p$, so it is irreducible, and $p=(a+bi)(a-bi)=(c+di)(c-di)$. By unique factorisation, $c+di$ is an associate of $a+bi$ or of $a-bi$, i.e.\ $c+di\in\{u(a\pm bi)\}$ with $u\in\{\pm1,\pm i\}$, which gives $\{|c|,|d|\}=\{|a|,|b|\}$.

(b) If $1000009$ were prime ($\equiv1\pmod4$), (a) would forbid the two representations. Euclid's algorithm in $\Z[i]$: $1000+3i=1\cdot(972+235i)+(28-232i)$, then $972+235i=4i(28-232i)+(44+123i)$, then $28-232i=(-2-i)(44+123i)+(-7+58i)$, and $44+123i=(2-i)(-7+58i)$ exactly (each quotient is the exact quotient rounded to a nearest Gaussian integer). So $g=-7+58i$, and $N(g)=3413$ divides $N(1000+3i)=1000009$: indeed $1000009=293\cdot3413$, with $293=17^2+2^2$ and $3413=58^2+7^2$.
\end{solution}

\begin{solution}{ex:x3y21}
\idea Factor the right side: $y^2+1=(y+i)(y-i)$. If the two factors are coprime, their product being a cube forces each to be a unit times a cube, and in $\Z[i]$ every unit is a cube. Then $y+i=(a+bi)^3$, and comparing imaginary parts gives a tiny equation. Coprimality needs $y$ even, which a congruence modulo $4$ provides.
\writeup If $y$ is odd, $y^2+1\equiv2\pmod4$, so $x^3$ is even but not divisible by $4$, impossible. So $y$ is even and $N(y+i)=y^2+1$ is odd. A common irreducible factor of $y+i$ and $y-i$ divides their difference $2i=(1+i)^2$, so it is $1+i$ up to a unit; but $N(1+i)=2$ does not divide the odd norm $y^2+1$. So they are coprime, and by unique factorisation $y+i=u(a+bi)^3$ with $u$ a unit. Every unit is a cube ($1=1^3$, $-1=(-1)^3$, $i=(-i)^3$, $-i=i^3$), so $y+i=(a+bi)^3=(a^3-3ab^2)+(3a^2b-b^3)i$. Then $b(3a^2-b^2)=1$, so $b=\pm1$ and $3a^2-1=\pm1$ with the same sign: $b=1$ gives $3a^2=2$, impossible; $b=-1$ gives $a=0$. So $y=a^3-3ab^2=0$ and $x=1$.
\end{solution}

\begin{solution}{ex:xxyy}
\idea The form $a^2-ab+b^2=N(a+b\omega)$ becomes $x^2+xy+y^2$ after $b\mapsto-b$. So $p$ is of this form iff $p$ is the norm of an Eisenstein integer, iff $p$ is not irreducible in $\Z[\omega]$; Theorem~\ref{thm:eisprimes} lists exactly when.
\writeup $x^2+xy+y^2=N(x-y\omega)$, and every $a+b\omega$ is $x-y\omega$ with $x=a$, $y=-b$. If $p=3$, $p=1+1+1$. If $p\equiv1\pmod3$, Theorem~\ref{thm:eisprimes} gives $\pi$ with $N(\pi)=p$. Conversely, if $p=N(z)$, then $p\ne2$ and $p\not\equiv2\pmod3$: otherwise $p$ is irreducible, and $p=z\bar z$ would be a product of two non-units (their norms are $p>1$). So $p=3$ or $p\equiv1\pmod3$.
\end{solution}

\begin{solution}{ex:y2x3m2}
\idea Same pattern as Exercise~\ref{ex:x3y21} in $\Z[\sqrt{-2}]$, which has unique factorisation by Proposition~\ref{prop:eucl}: $x^3=y^2+2=(y+\sqrt{-2})(y-\sqrt{-2})$. Check the factors are coprime (this needs $y$ odd), then $y+\sqrt{-2}=(a+b\sqrt{-2})^3$ since the units $\pm1$ are cubes, and compare coefficients of $\sqrt{-2}$.
\writeup If $y$ is even, then $x$ is even and $8\mid x^3=y^2+2\equiv2$ or $6\pmod8$: impossible. So $y$ is odd and $N(y\pm\sqrt{-2})=y^2+2$ is odd. A common factor of $y+\sqrt{-2}$ and $y-\sqrt{-2}$ divides $2\sqrt{-2}=-\sqrt{-2}^3$, so it is $\sqrt{-2}$ up to a unit, whose norm $2$ does not divide $y^2+2$. So the factors are coprime, and $y+\sqrt{-2}=\pm(a+b\sqrt{-2})^3=(a'+b'\sqrt{-2})^3$ (absorb the sign into the cube). Expanding, $(a+b\sqrt{-2})^3=a^3-6ab^2+(3a^2b-2b^3)\sqrt{-2}$, so $b(3a^2-2b^2)=1$: $b=\pm1$, and $3a^2-2=\pm1$ with the same sign. $b=1$ gives $a=\pm1$; $b=-1$ gives $3a^2=1$, impossible. Then $y=a^3-6a=\mp5$ and $x^3=25+2=27$, $x=3$. The solutions are $(3,\pm5)$.
\end{solution}

\section*{Sources and further reading}

\begin{itemize}
\item K.~Ireland and M.~Rosen, \emph{A Classical Introduction to Modern Number Theory}, 2nd ed., Springer, 1990, Chapters~1 and~9. Gaussian and Eisenstein integers, and cubic and biquadratic reciprocity.
\item D.~A.~Cox, \emph{Primes of the Form $x^2+ny^2$}, 3rd ed., AMS Chelsea, 2022, \S\S1 and~4. Where these rings lead: genus theory, class groups and class field theory.
\item H.~M.~Edwards, \emph{Fermat's Last Theorem: A Genetic Introduction to Algebraic Number Theory}, Springer, 1977, Chapters~2--4. Euler's proof for cubes, its gap, and Kummer's ideal numbers.
\item E.~Gethner, S.~Wagon and B.~Wick, A stroll through the Gaussian primes, \emph{American Mathematical Monthly} 105 (1998), 327--337. Can one walk to infinity on Gaussian primes with steps of bounded length? It is still an open problem.
\item F.~Lemmermeyer, \emph{Reciprocity Laws: From Euler to Eisenstein}, Springer, 2000, Chapters~6--7. The cubic and quartic reciprocity laws that live in these two rings.
\end{itemize}
